\documentclass[11pt]{article}
% Shared preamble for the printable mini lessons (1.4 onward).
%   \documentclass[11pt]{article}
%   % Shared preamble for the printable mini lessons (1.4 onward).
%   \documentclass[11pt]{article}
%   % Shared preamble for the printable mini lessons (1.4 onward).
%   \documentclass[11pt]{article}
%   \input{../preamble}
% Build the key with:  pdflatex -jobname=<name>_KEY "\def\iskey{}\input{<name>}"
% Figures shared with the web page live in figures/Unit_X/<name>.tex and are
% pulled in with \fig{Unit_X/<name>}; the standalone package drops their own
% preamble, so their macros come from figures/preamble.tex, loaded here.
\input{../../figures/preamble}
\usepackage[margin=0.75in]{geometry}
\usepackage{array}
\usepackage{enumitem}
\usepackage{amsmath}
\usepackage[helvet]{sfmath}  % math in the same sans font as the text
\usepackage{standalone}
\pagestyle{empty}
\setlength{\parindent}{0pt}
\setlist[enumerate,1]{leftmargin=1.5em,itemsep=1em}
\setlist[enumerate,2]{label=(\alph*),leftmargin=2em,itemsep=2.2em}

\newif\ifkey
\ifdefined\iskey \keytrue \fi
\colorlet{keyred}{red!75!black}
\newcommand{\ans}[1]{\ifkey\leavevmode{\color{keyred}#1}\fi}
% vertical work space on the student copy, less on the key
\newcommand{\work}[1]{\ifkey\vspace{0.6em}\else\vspace{#1}\fi}
\newcommand{\fig}[1]{\input{../../figures/#1}}
% Title line: \lessontitle{Quartiles, IQR, and Box Plots}
\newcommand{\lessontitle}[1]{{\Large\bfseries #1\ifkey\ \textcolor{keyred}{KEY}\fi}\hfill Name \underline{\hspace{2in}}\par\vspace{0.6em}}
\newcommand{\partner}{\par\vspace{0.4em}Work with a partner. Compare answers when you finish each problem.\par}

% Build the key with:  pdflatex -jobname=<name>_KEY "\def\iskey{}\input{<name>}"
% Figures shared with the web page live in figures/Unit_X/<name>.tex and are
% pulled in with \fig{Unit_X/<name>}; the standalone package drops their own
% preamble, so their macros come from figures/preamble.tex, loaded here.
% Shared preamble for web figures.  Each figure is a standalone .tex that
% \input's this file; scripts/make-figures.sh turns it into an SVG.
\usepackage[T1]{fontenc}
\usepackage{helvet}
\renewcommand{\familydefault}{\sfdefault}
\usepackage{tikz}
\usetikzlibrary{decorations.pathreplacing}
% axis numbers in the text font, not math italic
\pgfkeys{/pgf/number format/.cd,assume math mode=true,1000 sep={}}

% dot plot: \dotplot{min}{max}{value/count,...}{axis label}
\newcommand{\dotplot}[4]{%
\begin{tikzpicture}[x=0.9cm,y=0.33cm]
  \draw[thick] (#1-0.5,0) -- (#2+0.5,0);
  \foreach \t in {#1,...,#2} {\draw (\t,0.25) -- (\t,-0.25) node[below] {\t};}
  \foreach \v/\c in {#3} {\ifnum\c>0 \foreach \k in {1,...,\c} {\fill (\v,\k*1.15) circle (0.16cm);}\fi}
  \node at ({(#1+#2)/2},-2.3) {#4};
\end{tikzpicture}}

% shape sketch histogram (12 bins): \shapehist{heights}{label}
\newcommand{\shapehist}[2]{%
  \foreach \h [count=\i from 0] in {#1} \draw[fill=gray!30] (\i,0) rectangle (\i+1,\h);
  \draw (-0.3,0) -- (12.3,0);
  \node[font=\bfseries] at (6,-2.6) {#2};}
\newcommand{\meanmark}[1]{\fill (#1,0) -- ++(-0.3,-0.75) -- ++(0.6,0) -- cycle;}
\newcommand{\medmark}[2]{\draw[very thick,dashed] (#1,0) -- (#1,#2);}

% number line with minor ticks: \nline{min}{max}{label step}{axis label}
\newcommand{\nline}[4]{%
  \draw[thick] (#1-0.5,0) -- (#2+0.5,0);
  \foreach \t in {#1,...,#2} \draw (\t,0.12) -- (\t,-0.12);
  \pgfmathtruncatemacro{\nxt}{#1+#3}
  \foreach \t in {#1,\nxt,...,#2} \draw (\t,0.2) -- (\t,-0.2) node[below] {\t};
  \node at ({(#1+#2)/2},-0.95) {#4};}
% box plot above the line: \bplot{min}{Q1}{med}{Q3}{max}
\newcommand{\bplot}[5]{%
  \draw[thick] (#1,1.0) -- (#2,1.0);
  \draw[thick] (#4,1.0) -- (#5,1.0);
  \draw[thick] (#1,0.75) -- (#1,1.25);
  \draw[thick] (#5,0.75) -- (#5,1.25);
  \draw[thick,fill=gray!20] (#2,0.55) rectangle (#4,1.45);
  \draw[very thick] (#3,0.55) -- (#3,1.45);}

% ---- added for 1.4 onward ----
% axis on y=0: \axis{min}{max}{label step}{tick step}{axis label}
\newcommand{\axis}[5]{%
  \draw[thick] (#1,0) -- (#2,0);
  \pgfmathsetmacro{\axT}{#1+#4}\foreach \t in {#1,\axT,...,#2} \draw (\t,0.1) -- (\t,-0.1);
  \pgfmathsetmacro{\axL}{#1+#3}\foreach \t in {#1,\axL,...,#2}
    \draw (\t,0.18) -- (\t,-0.18) node[below] {\pgfmathprintnumber{\t}};
  \node at ({(#1+#2)/2},-0.95) {#5};}
% box plot centred at height y: \bplotat{y}{min}{Q1}{med}{Q3}{max}
\newcommand{\bplotat}[6]{%
  \draw[thick] (#2,#1) -- (#3,#1);
  \draw[thick] (#5,#1) -- (#6,#1);
  \draw[thick] (#2,#1-0.22) -- (#2,#1+0.22);
  \draw[thick] (#6,#1-0.22) -- (#6,#1+0.22);
  \draw[thick,fill=gray!20] (#3,#1-0.42) rectangle (#5,#1+0.42);
  \draw[very thick] (#4,#1-0.42) -- (#4,#1+0.42);}
% an outlier dot: \outlier{x}{y}
\newcommand{\outlier}[2]{\fill (#1,#2) circle (2.4pt);}

\usepackage[margin=0.75in]{geometry}
\usepackage{array}
\usepackage{enumitem}
\usepackage{amsmath}
\usepackage[helvet]{sfmath}  % math in the same sans font as the text
\usepackage{standalone}
\pagestyle{empty}
\setlength{\parindent}{0pt}
\setlist[enumerate,1]{leftmargin=1.5em,itemsep=1em}
\setlist[enumerate,2]{label=(\alph*),leftmargin=2em,itemsep=2.2em}

\newif\ifkey
\ifdefined\iskey \keytrue \fi
\colorlet{keyred}{red!75!black}
\newcommand{\ans}[1]{\ifkey\leavevmode{\color{keyred}#1}\fi}
% vertical work space on the student copy, less on the key
\newcommand{\work}[1]{\ifkey\vspace{0.6em}\else\vspace{#1}\fi}
\newcommand{\fig}[1]{\input{../../figures/#1}}
% Title line: \lessontitle{Quartiles, IQR, and Box Plots}
\newcommand{\lessontitle}[1]{{\Large\bfseries #1\ifkey\ \textcolor{keyred}{KEY}\fi}\hfill Name \underline{\hspace{2in}}\par\vspace{0.6em}}
\newcommand{\partner}{\par\vspace{0.4em}Work with a partner. Compare answers when you finish each problem.\par}

% Build the key with:  pdflatex -jobname=<name>_KEY "\def\iskey{}\input{<name>}"
% Figures shared with the web page live in figures/Unit_X/<name>.tex and are
% pulled in with \fig{Unit_X/<name>}; the standalone package drops their own
% preamble, so their macros come from figures/preamble.tex, loaded here.
% Shared preamble for web figures.  Each figure is a standalone .tex that
% \input's this file; scripts/make-figures.sh turns it into an SVG.
\usepackage[T1]{fontenc}
\usepackage{helvet}
\renewcommand{\familydefault}{\sfdefault}
\usepackage{tikz}
\usetikzlibrary{decorations.pathreplacing}
% axis numbers in the text font, not math italic
\pgfkeys{/pgf/number format/.cd,assume math mode=true,1000 sep={}}

% dot plot: \dotplot{min}{max}{value/count,...}{axis label}
\newcommand{\dotplot}[4]{%
\begin{tikzpicture}[x=0.9cm,y=0.33cm]
  \draw[thick] (#1-0.5,0) -- (#2+0.5,0);
  \foreach \t in {#1,...,#2} {\draw (\t,0.25) -- (\t,-0.25) node[below] {\t};}
  \foreach \v/\c in {#3} {\ifnum\c>0 \foreach \k in {1,...,\c} {\fill (\v,\k*1.15) circle (0.16cm);}\fi}
  \node at ({(#1+#2)/2},-2.3) {#4};
\end{tikzpicture}}

% shape sketch histogram (12 bins): \shapehist{heights}{label}
\newcommand{\shapehist}[2]{%
  \foreach \h [count=\i from 0] in {#1} \draw[fill=gray!30] (\i,0) rectangle (\i+1,\h);
  \draw (-0.3,0) -- (12.3,0);
  \node[font=\bfseries] at (6,-2.6) {#2};}
\newcommand{\meanmark}[1]{\fill (#1,0) -- ++(-0.3,-0.75) -- ++(0.6,0) -- cycle;}
\newcommand{\medmark}[2]{\draw[very thick,dashed] (#1,0) -- (#1,#2);}

% number line with minor ticks: \nline{min}{max}{label step}{axis label}
\newcommand{\nline}[4]{%
  \draw[thick] (#1-0.5,0) -- (#2+0.5,0);
  \foreach \t in {#1,...,#2} \draw (\t,0.12) -- (\t,-0.12);
  \pgfmathtruncatemacro{\nxt}{#1+#3}
  \foreach \t in {#1,\nxt,...,#2} \draw (\t,0.2) -- (\t,-0.2) node[below] {\t};
  \node at ({(#1+#2)/2},-0.95) {#4};}
% box plot above the line: \bplot{min}{Q1}{med}{Q3}{max}
\newcommand{\bplot}[5]{%
  \draw[thick] (#1,1.0) -- (#2,1.0);
  \draw[thick] (#4,1.0) -- (#5,1.0);
  \draw[thick] (#1,0.75) -- (#1,1.25);
  \draw[thick] (#5,0.75) -- (#5,1.25);
  \draw[thick,fill=gray!20] (#2,0.55) rectangle (#4,1.45);
  \draw[very thick] (#3,0.55) -- (#3,1.45);}

% ---- added for 1.4 onward ----
% axis on y=0: \axis{min}{max}{label step}{tick step}{axis label}
\newcommand{\axis}[5]{%
  \draw[thick] (#1,0) -- (#2,0);
  \pgfmathsetmacro{\axT}{#1+#4}\foreach \t in {#1,\axT,...,#2} \draw (\t,0.1) -- (\t,-0.1);
  \pgfmathsetmacro{\axL}{#1+#3}\foreach \t in {#1,\axL,...,#2}
    \draw (\t,0.18) -- (\t,-0.18) node[below] {\pgfmathprintnumber{\t}};
  \node at ({(#1+#2)/2},-0.95) {#5};}
% box plot centred at height y: \bplotat{y}{min}{Q1}{med}{Q3}{max}
\newcommand{\bplotat}[6]{%
  \draw[thick] (#2,#1) -- (#3,#1);
  \draw[thick] (#5,#1) -- (#6,#1);
  \draw[thick] (#2,#1-0.22) -- (#2,#1+0.22);
  \draw[thick] (#6,#1-0.22) -- (#6,#1+0.22);
  \draw[thick,fill=gray!20] (#3,#1-0.42) rectangle (#5,#1+0.42);
  \draw[very thick] (#4,#1-0.42) -- (#4,#1+0.42);}
% an outlier dot: \outlier{x}{y}
\newcommand{\outlier}[2]{\fill (#1,#2) circle (2.4pt);}

\usepackage[margin=0.75in]{geometry}
\usepackage{array}
\usepackage{enumitem}
\usepackage{amsmath}
\usepackage[helvet]{sfmath}  % math in the same sans font as the text
\usepackage{standalone}
\pagestyle{empty}
\setlength{\parindent}{0pt}
\setlist[enumerate,1]{leftmargin=1.5em,itemsep=1em}
\setlist[enumerate,2]{label=(\alph*),leftmargin=2em,itemsep=2.2em}

\newif\ifkey
\ifdefined\iskey \keytrue \fi
\colorlet{keyred}{red!75!black}
\newcommand{\ans}[1]{\ifkey\leavevmode{\color{keyred}#1}\fi}
% vertical work space on the student copy, less on the key
\newcommand{\work}[1]{\ifkey\vspace{0.6em}\else\vspace{#1}\fi}
\newcommand{\fig}[1]{\input{../../figures/#1}}
% Title line: \lessontitle{Quartiles, IQR, and Box Plots}
\newcommand{\lessontitle}[1]{{\Large\bfseries #1\ifkey\ \textcolor{keyred}{KEY}\fi}\hfill Name \underline{\hspace{2in}}\par\vspace{0.6em}}
\newcommand{\partner}{\par\vspace{0.4em}Work with a partner. Compare answers when you finish each problem.\par}


\begin{document}

\lessontitle{Reading Box Plots and Outliers}

\textbf{Reading a box plot.} Each of the 4 parts holds about \textbf{25\%} of the data.

\begin{center}\fig{Unit_1/1.4_box_plot_scores}\end{center}

\begin{tabular}{@{}ll}
The box holds the middle \textbf{50\%}: & half of the students scored from 68 to 84.\\
Above Q1 is \textbf{75\%}: & about 75\% of the students scored more than 68.\\
\end{tabular}

\vspace{0.6em}
\textbf{Comparing box plots.} Use the same number line. Compare the \textbf{centers} and the \textbf{spreads}.

\begin{center}\fig{Unit_1/1.4_compare_classes}\end{center}

\begin{tabular}{@{}ll}
Center (median): & Afternoon 78, Morning 72. The afternoon class scored higher.\\
Spread (IQR): & Afternoon $82 - 70 = 12$, Morning $80 - 65 = 15$. Afternoon scores are closer together.\\
\end{tabular}

\vspace{0.8em}
\textbf{Outliers.} An \textbf{outlier} is a value far away from the rest of the data.

\vspace{0.3em}
\renewcommand{\arraystretch}{1.3}
\begin{tabular}{|>{\raggedright\arraybackslash}p{3.1in}|>{\raggedright\arraybackslash}p{3.4in}|}
\hline
\textbf{The 1.5 $\times$ IQR rule} & \textbf{Example:} 3, 8, 9, 10, 11, 12, 13, 14, 30 \\
1. Find Q1, Q3, and the IQR. &
Q1 = 8.5, \; Q3 = 13.5, \; IQR = 5 \\
2. Lower fence = Q1 $-$ 1.5 $\times$ IQR &
$8.5 - 1.5 \times 5 = 8.5 - 7.5 = 1$ \\
3. Upper fence = Q3 + 1.5 $\times$ IQR &
$13.5 + 7.5 = 21$ \\
4. A value below the lower fence or above the upper fence is an outlier. &
30 is above 21, so \textbf{30 is an outlier}. 3 is not below 1, so 3 is not an outlier. \\
\hline
\end{tabular}

\begin{center}\fig{Unit_1/1.4_outlier_example}\end{center}

On a box plot, an outlier is a \textbf{dot}. The whisker stops at the last value that is not an outlier (here, 14).

\partner

% ---------------- PROBLEMS ----------------
\newpage
\begin{enumerate}

\item This box plot shows how many minutes some students exercise each day.
\begin{center}\fig{Unit_1/1.4_box_plot_exercise}\end{center}
\begin{enumerate}[itemsep=1.6em]
  \item Write the five-number summary. \ans{\quad Min 10, Q1 20, Median 35, Q3 45, Max 70}
  \item Find the IQR. \ans{\quad $45 - 20 = 25$ minutes}
  \item Find the range. \ans{\quad $70 - 10 = 60$ minutes}
  \item About what percent of the students exercise more than 20 minutes? \ans{\quad 75\%}
  \item About what percent of the students exercise from 35 to 70 minutes? \ans{\quad 50\%}
\end{enumerate}
\work{1.5em}

\item These box plots show how many hours students in two classes use their phones each day.
\begin{center}\fig{Unit_1/1.4_box_plots_phone}\end{center}
\begin{enumerate}[itemsep=1.6em]
  \item Which class has the higher median? \ans{\quad Class B (5 hours vs 4 hours)}
  \item Find the IQR of each class. \ans{\quad A: $6 - 3 = 3$ hours. \; B: $5.5 - 4 = 1.5$ hours.}
  \item Which class is more \textbf{consistent} (the values are closer together)? Explain.

  \ans{Class B. Its IQR and its range are both smaller.}
  \item About what percent of Class B uses the phone more than 4 hours a day? \ans{\quad 75\%}
\end{enumerate}

\newpage
\item Nine students recorded how many minutes they ride the bus to school.
\begin{center}12,\; 15,\; 16,\; 18,\; 19,\; 20,\; 21,\; 22,\; 45\end{center}
\begin{enumerate}[itemsep=1em]
  \item Find Q1, Q3, and the IQR.

  \ans{Median 19. Lower half 12, 15, 16, 18: Q1 = 15.5. Upper half 20, 21, 22, 45: Q3 = 21.5. IQR = 6.}
  \work{4em}
  \item Find the lower fence and the upper fence.

  \ans{$15.5 - 1.5 \times 6 = 15.5 - 9 = 6.5$. \quad $21.5 + 9 = 30.5$}
  \work{4em}
  \item Are there any outliers? Explain.

  \ans{Yes. 45 is above the upper fence of 30.5. No value is below 6.5.}
\end{enumerate}
\work{3em}

\item Ten students timed how many minutes they took to finish a quiz.
\begin{center}23,\; 45,\; 20,\; 27,\; 2,\; 24,\; 21,\; 28,\; 18,\; 25\end{center}
\begin{enumerate}[itemsep=1em]
  \item Find Q1, the median, Q3, and the IQR.

  \ans{2, 18, 20, 21, \fbox{23}, \fbox{24}, 25, 27, 28, 45. \quad Median 23.5\\
  Lower half 2, 18, 20, 21, 23: Q1 = 20. \quad Upper half 24, 25, 27, 28, 45: Q3 = 27. \quad IQR = 7}
  \work{5em}
  \item Find the fences. Which values are outliers?

  \ans{$20 - 10.5 = 9.5$ \; and \; $27 + 10.5 = 37.5$. \quad 2 and 45 are outliers.}
  \work{4em}
  \item Make a box plot. Show each outlier as a dot.
  \begin{center}
  \begin{tikzpicture}[x=0.3cm,y=1cm]
    \axis{0}{50}{5}{1}{Minutes to finish the quiz}
    \ifkey\begin{scope}[keyred]\bplotat{1}{18}{20}{23.5}{27}{28}\outlier{2}{1}\outlier{45}{1}\end{scope}\else\path (0,1.6);\fi
  \end{tikzpicture}
  \end{center}
\end{enumerate}

\newpage
\item This box plot shows the points scored by players on a basketball team.
\begin{center}\fig{Unit_1/1.4_box_plot_points}\end{center}
\begin{enumerate}[itemsep=1.6em]
  \item What is the maximum? \ans{\quad 40 points (the dot)}
  \item Find the IQR. \ans{\quad $14 - 6 = 8$ points}
  \item Use the 1.5 $\times$ IQR rule to show that 40 is an outlier.

  \ans{Upper fence $= 14 + 1.5 \times 8 = 14 + 12 = 26$. 40 is above 26.}
  \work{3em}
  \item Why does the whisker stop at 22?

  \ans{22 is the largest value that is not an outlier.}
\end{enumerate}
\work{2em}

\item Look at the exercise box plot in problem 1. Is each statement true or false? Explain.
\begin{enumerate}[itemsep=2.4em]
  \item About half of the students exercise from 20 to 45 minutes each day.

  \ans{True. The box holds the middle 50\%.}
  \item The right whisker is the longest part, so most students are in that part.

  \ans{False. Each part has about 25\% of the students. A long part means those values are spread out.}
  \item The mean is 35 minutes.

  \ans{False (or we cannot tell). 35 is the median. A box plot does not show the mean.}
  \item At least one student exercises exactly 70 minutes.

  \ans{True. The maximum is a value in the data.}
\end{enumerate}

\end{enumerate}

\end{document}
